% ============================================================================
% MINITEST MUSTERLÖSUNG: Mi instituto - Objetos de clase & hay
% Klasse 7 | Spanisch | 25 BE | Papierversion
% ============================================================================

\documentclass[11pt,a4paper]{article}

\usepackage[utf8]{inputenc}
\usepackage[T1]{fontenc}
\usepackage[ngerman]{babel}
\usepackage{geometry}
\usepackage{fancyhdr}
\usepackage{xcolor}
\usepackage{enumitem}
\usepackage{tabularx}
\usepackage{array}
\usepackage{tcolorbox}
\usepackage{amssymb}

% ============================================================================
% LAYOUT
% ============================================================================
\geometry{left=2cm, right=2cm, top=2cm, bottom=2cm}
\definecolor{fach}{HTML}{c41e3a}
\definecolor{grau}{HTML}{666666}
\definecolor{loesung}{HTML}{c62828}

\pagestyle{fancy}
\fancyhf{}
\fancyhead[L]{\small\textcolor{grau}{Spanisch Klasse 7 -- Minitest}}
\fancyhead[R]{\small\textcolor{loesung}{\textbf{MUSTERLÖSUNG}} MT-01-ML}
\fancyfoot[C]{\small\thepage}
\renewcommand{\headrulewidth}{0.4pt}

% Befehle
\newcommand{\luecke}[1]{\underline{\hspace{#1}}}
\newcommand{\lsg}[1]{\textcolor{loesung}{\textbf{#1}}}
\newcommand{\punktebox}[1]{\hfill\fbox{\small \_/#1}}
\newcommand{\kreisLeer}{$\bigcirc$}
\newcommand{\kreisVoll}{\textcolor{loesung}{$\CIRCLE$}}

% ============================================================================
% DOKUMENT
% ============================================================================
\begin{document}

% ----------------------------------------------------------------------------
% KOPFBEREICH
% ----------------------------------------------------------------------------
\begin{center}
{\Large\textcolor{fach}{\textbf{Minitest: Objetos de clase \& hay}}}\\[0.3cm]
{\normalsize Unidad 2 | 25 BE | ca. 15 Minuten}\\[0.2cm]
{\large\textcolor{loesung}{\textbf{--- MUSTERLÖSUNG ---}}}
\end{center}

\vspace{0.3cm}
\noindent\textbf{Name:} \luecke{5cm} \hfill \textbf{Datum:} \luecke{3cm} \hfill \textbf{Klasse:} \luecke{2cm}

\vspace{0.3cm}
\hrule
\vspace{0.5cm}

% ----------------------------------------------------------------------------
% AUFGABE 1: Multiple Choice Vokabeln
% ----------------------------------------------------------------------------
\noindent\textbf{Aufgabe 1: Vokabeln} \punktebox{8}

\noindent Kreuze die richtige Übersetzung an.

\vspace{0.3cm}

\begin{tabular}{p{3cm}p{4cm}p{4cm}p{4cm}}
\textbf{1.1 la mesa} & \kreisLeer{} der Stuhl & \kreisVoll{} \lsg{der Tisch} & \kreisLeer{} die Tafel \\[0.5em]
\textbf{1.2 la silla} & \kreisVoll{} \lsg{der Stuhl} & \kreisLeer{} das Fenster & \kreisLeer{} die Tür \\[0.5em]
\textbf{1.3 el libro} & \kreisLeer{} das Heft & \kreisLeer{} der Kuli & \kreisVoll{} \lsg{das Buch} \\[0.5em]
\textbf{1.4 la pizarra} & \kreisVoll{} \lsg{die Tafel} & \kreisLeer{} die Wand & \kreisLeer{} das Fenster \\
\end{tabular}

\textcolor{loesung}{\textit{(Je 2 BE pro richtige Antwort = 8 BE)}}

\vspace{0.8cm}

% ----------------------------------------------------------------------------
% AUFGABE 2: Artikel
% ----------------------------------------------------------------------------
\noindent\textbf{Aufgabe 2: Artikel} \punktebox{6}

\noindent Ergänze \textbf{el} oder \textbf{la}.

\vspace{0.3cm}

\begin{tabular}{p{6cm}p{6cm}}
2.1 \lsg{la} mesa & 2.2 \lsg{el} libro \\[0.8em]
2.3 \lsg{el} cuaderno & 2.4 \lsg{la} ventana \\[0.8em]
2.5 \lsg{el} mapa \textit{(Achtung: maskulin!)} & 2.6 \lsg{la} puerta \\
\end{tabular}

\textcolor{loesung}{\textit{(Je 1 BE pro richtige Antwort = 6 BE)}}

\vspace{0.8cm}

% ----------------------------------------------------------------------------
% AUFGABE 3: hay-Lücken
% ----------------------------------------------------------------------------
\noindent\textbf{Aufgabe 3: hay} \punktebox{6}

\noindent Ergänze \textbf{hay}, \textbf{un}, \textbf{una} oder \textbf{muchos/muchas}.

\vspace{0.3cm}

\begin{enumerate}[label=3.\arabic*]
    \item En mi clase \lsg{hay} una pizarra.
    \item \lsg{Hay} veinticinco sillas.
    \item En la pared hay \lsg{un} mapa.
    \item En mi mochila hay \lsg{un} libro.
    \item Hay \lsg{muchos} cuadernos. \textit{(viele)}
    \item ¿Qué \lsg{hay} en la clase?
\end{enumerate}

\textcolor{loesung}{\textit{(Je 1 BE pro richtige Antwort = 6 BE)}}

\vspace{0.8cm}

% ----------------------------------------------------------------------------
% AUFGABE 4: Fehlerkorrektur
% ----------------------------------------------------------------------------
\noindent\textbf{Aufgabe 4: Fehlerkorrektur} \punktebox{5}

\noindent Welcher Satz ist \textbf{falsch}? Kreuze an.

\vspace{0.3cm}

\textbf{4.1}
\begin{itemize}[label=$\bigcirc$, nosep]
    \item Hay una mesa.
    \item[\kreisVoll] \lsg{Hay el libro.} \textcolor{loesung}{\textit{(FALSCH -- kein best. Artikel nach hay!)}}
    \item Hay muchos cuadernos.
\end{itemize}

\vspace{0.3cm}

\textbf{4.2}
\begin{itemize}[label=$\bigcirc$, nosep]
    \item Hay cinco sillas.
    \item[\kreisVoll] \lsg{Hay la ventana.} \textcolor{loesung}{\textit{(FALSCH -- kein best. Artikel nach hay!)}}
    \item En mi clase hay una pizarra.
\end{itemize}

\textcolor{loesung}{\textit{(Je 2 BE pro richtige Antwort = 4 BE)}}

\vspace{0.5cm}

\noindent\textbf{4.3 Korrigiere:} *Hay el libro $\rightarrow$ Hay \lsg{un} libro.

\textcolor{loesung}{\textit{(1 BE)}}

\vspace{1cm}

% ----------------------------------------------------------------------------
% PUNKTEZEILE
% ----------------------------------------------------------------------------
\begin{flushright}
\textbf{Punkte:} \luecke{1cm} / 25
\end{flushright}

\vspace{0.5cm}

% ----------------------------------------------------------------------------
% NOTENTABELLE
% ----------------------------------------------------------------------------
\begin{tcolorbox}[colback=white, colframe=grau, title=\textbf{Bewertung (14-NP-Skala)}]
\begin{center}
\small
\begin{tabular}{|c|c|c|c|c|c|c|c|c|c|c|c|c|c|c|}
\hline
\textbf{NP} & 14 & 13 & 12 & 11 & 10 & 9 & 8 & 7 & 6 & 5 & 4 & 3 & 2 & 1 \\
\hline
\textbf{BE} & 25 & 24 & 23 & 22 & 20 & 18 & 17 & 15 & 13 & 12 & 10 & 8 & 7 & 5 \\
\hline
\textbf{\%} & 100 & 96 & 92 & 88 & 80 & 72 & 68 & 60 & 52 & 48 & 40 & 32 & 28 & 20 \\
\hline
\end{tabular}
\end{center}
\end{tcolorbox}

\vspace{0.5cm}

% ----------------------------------------------------------------------------
% ERWARTUNGSHORIZONT KURZFORM
% ----------------------------------------------------------------------------
\begin{tcolorbox}[colback=white, colframe=loesung, title=\textbf{Erwartungshorizont (Kurzform)}]
\begin{tabular}{|l|l|l|}
\hline
\textbf{Aufgabe} & \textbf{Lösung} & \textbf{BE} \\
\hline
1.1 & b) der Tisch & 2 \\
1.2 & a) der Stuhl & 2 \\
1.3 & c) das Buch & 2 \\
1.4 & a) die Tafel & 2 \\
\hline
2.1--2.6 & la, el, el, la, el, la & 6 \\
\hline
3.1--3.6 & hay, Hay, un, un, muchos, hay & 6 \\
\hline
4.1 & b) Hay el libro & 2 \\
4.2 & b) Hay la ventana & 2 \\
4.3 & un & 1 \\
\hline
\textbf{Gesamt} & & \textbf{25} \\
\hline
\end{tabular}
\end{tcolorbox}

\end{document}
